\chapter{Source Texts and Their Coverage}
\label{ch:sources}

The course requires that any textbook consulted beyond the prescribed one be
identified in full, and that the portion of its table of contents relied upon be
reproduced so that the correspondence with this paper is verifiable.
This chapter
does that.
It also states, for each source, which chapter of this paper draws on
it and for what.

Sources are grouped into three tiers.
Tier~1 texts are read in full and carry the
technical weight of the paper.
Tier~2 texts are consulted for specific treatments
and cross-checks.
Tier~3 comprises the vendor architecture manuals, which are the
authoritative statement of what each shipping machine actually promises and are
therefore the primary source wherever this paper makes a claim about a particular
instruction set.

\section{Tier 1: Primary Texts}

\subsection{Sarangi, \emph{Basic Computer Architecture} (prescribed)}

\begin{table}[H]
\centering
\small
\begin{tabular}{@{}ll@{}}
\toprule
Author      & Smruti R. Sarangi \\
Title       & Basic Computer Architecture \\
Version     & 3.09 (8 October 2025) \\
Publisher   & White Falcon Publishing \\
Licence     & CC BY-ND 4.0; PDF freely available \\
URL         & \url{https://srsarangi.github.io/archbook/archbook.pdf} \\
\bottomrule
\end{tabular}
\end{table}

\noindent Relevant portion of the table of contents:

\begin{table}[H]
\centering
\small
\begin{tabular}{@{}llr@{}}
\toprule
\textbf{\S} & \textbf{Title} & \textbf{Page} \\
\midrule
\multicolumn{3}{@{}l}{\textit{Chapter 10 --- Principles of Pipelining}} \\
10.4      & Pipeline Hazards                              & 420 \\
10.9      & Performance Metrics                           & 463 \\
10.11.4   & Out-of-Order Pipelines                        & 492 \\
\addlinespace
\multicolumn{3}{@{}l}{\textit{Chapter 11 --- The Memory System}} \\
11.1      & Overview (locality, hierarchy, cache blocks)  & 506 \\
11.2      & Caches                                        & 517 \\
11.3      & The Memory System                             & 532 \\
\addlinespace
\multicolumn{3}{@{}l}{\textit{Chapter 12 --- Multiprocessor Systems}} \\
12.2.2    & Shared Memory vs Message Passing              & 571 \\
12.2.3    & Amdahl's Law                                  & 576 \\
12.3      & Design Space of Multiprocessors               & 578 \\
12.4      & MIMD Multiprocessors                          & 579 \\
12.4.1    & Logical Point of View                         & 579 \\
12.4.2    & Coherence                                     & 581 \\
12.4.3    & \textbf{Memory Consistency}                   & 583 \\
12.4.4    & Physical View of Memory                       & 593 \\
12.4.5    & Shared Caches                                 & 596 \\
12.4.6    & Coherent Private Caches                       & 596 \\
12.4.7    & Implementing a Memory Consistency Model*      & 603 \\
12.4.8    & Multithreaded Processors                      & 608 \\
12.6      & Interconnection Networks                      & 620 \\
\addlinespace
\multicolumn{3}{@{}l}{\textit{Appendix A --- Case Studies of Real Processors}} \\
A.1.2     & ARM Cortex-A8                                 & 734 \\
A.1.3     & ARM Cortex-A15                                & 736 \\
A.3.2     & Intel Sandy Bridge                            & 746 \\
\bottomrule
\end{tabular}
\caption{Sections of the prescribed text spanned by this paper. Sections marked
$*$ are designated advanced in the text.}
\label{tab:toc-archbook}
\end{table}

\noindent \textbf{Use in this paper.} \S12.4.2 and \S12.4.3 are the sections this
paper extends; they supply the coherence/consistency distinction developed in
Chapter~3 and the sequential-consistency treatment developed in Chapter~4.
\S12.4.7 is the direct antecedent of the implementation discussion. \S11.2 and
\S10.11.4 supply the microarchitectural background of Chapter~2, \S10.9 the
performance-metric framework used in Chapter~6, and \S12.6 the interconnect
parameters of the simulator.

\subsection{Sarangi, \emph{Next-Gen Computer Architecture}}

\begin{table}[H]
\centering
\small
\begin{tabular}{@{}ll@{}}
\toprule
Author      & Smruti R. Sarangi \\
Title       & Next-Gen Computer Architecture: Till the End of Silicon \\
Publisher   & White Falcon Publishing, October 2023 \\
ISBN        & 8119510143 \\
Note        & First edition published by McGraw-Hill (2021) as
              \emph{Advanced Computer Architecture}, ISBN 9789390727490 \\
Licence     & CC BY-ND 4.0; PDF freely available \\
URL         & \url{https://www.cse.iitd.ac.in/~srsarangi/advbook/} \\
\bottomrule
\end{tabular}
\end{table}

This is the companion volume to the prescribed text, aimed at final-year
undergraduates, and it is the single most important source for this paper after
the prescribed book.
Chapter~9 is devoted entirely to the subject and runs to
approximately 118 pages.

\begin{table}[H]
\centering
\small
\begin{tabular}{@{}lp{7.4cm}l@{}}
\toprule
\textbf{\S} & \textbf{Title} & \textbf{Used in} \\
\midrule
\multicolumn{3}{@{}l}{\textit{Chapter 9 --- Multicore Systems: Coherence,
  Consistency and Transactional Memory}} \\
9.1 & Parallel programming and hardware threads          & Ch.~2 \\
9.2 & Theoretical foundations                            & Ch.~4 \\
9.3 & Sequential consistency, PLSC, and coherence        & Ch.~3, Ch.~4 \\
9.4 & \textbf{Execution witnesses, access graphs, causal graphs} & Ch.~4 \\
9.5 & Cache coherence: snoopy and directory protocols    & Ch.~3 \\
9.6 & Advanced directory protocols and atomic operations & Ch.~3, Ch.~8 \\
9.7 & Memory models and data races                       & Ch.~4, Ch.~6 \\
9.8 & Methods to detect races                            & Ch.~10 \\
9.9 & Transactional memory                               & Ch.~11 \\
\addlinespace
\multicolumn{3}{@{}l}{\textit{Supporting chapters}} \\
4.3 & Load store queue                                   & Ch.~2 \\
5.1 & Aggressive speculation                             & Ch.~9 \\
7   & Caches                                             & Ch.~2 \\
8   & On-chip network                                    & Ch.~8 \\
App.~B & Tejas Architectural Simulator                   & Ch.~8 \\
\bottomrule
\end{tabular}
\caption{Sections of \emph{Next-Gen Computer Architecture} consulted. Page
numbers are omitted because they differ between the McGraw-Hill and White Falcon
printings and between PDF versions; section numbering is stable.}
\label{tab:toc-nextgen}
\end{table}

\noindent \textbf{Use in this paper.} \S9.4 supplies the formalism of Chapter~4.
The execution-witness and access-graph notation used throughout this paper is
Sarangi's, chosen in preference to the alternative formulations in the literature
so that the notation is continuous with the course.
\S9.3's treatment of
per-location sequential consistency is the source of the PLSC characterisation of
coherence in \S\ref{sec:invariants}.
\S9.7 and \S9.8 underpin the data-race
material of \S\ref{sec:drf} and the verification design of Chapter~10.
Appendix~B documents the Tejas simulator, evaluated as an alternative to the
purpose-built simulator of Chapter~8.

\subsection{Nagarajan, Sorin, Hill and Wood, \emph{A Primer on Memory
Consistency and Cache Coherence}}

\begin{table}[H]
\centering
\small
\begin{tabular}{@{}ll@{}}
\toprule
Authors     & Vijay Nagarajan, Daniel J. Sorin, Mark D. Hill, David A. Wood \\
Title       & A Primer on Memory Consistency and Cache Coherence \\
Edition     & Second edition, 2020 \\
Series      & Synthesis Lectures on Computer Architecture \\
Publisher   & Springer Nature, Cham \\
Access      & Open access via the OAPEN Library \\
\bottomrule
\end{tabular}
\end{table}

\begin{table}[H]
\centering
\small
\begin{tabular}{@{}rp{8.2cm}l@{}}
\toprule
\textbf{Ch.} & \textbf{Title} & \textbf{Used in} \\
\midrule
1  & Introduction to Consistency and Coherence            & Ch.~1 \\
2  & Coherence Basics                                     & Ch.~3 \\
3  & Memory Consistency Motivation and Sequential Consistency & Ch.~1, Ch.~4 \\
4  & Total Store Order and the x86 Memory Model           & Ch.~4, Ch.~7 \\
5  & Relaxed Memory Consistency                           & Ch.~4 \\
6  & Coherence Protocols                                  & Ch.~3 \\
7  & Snooping Coherence Protocols                         & Ch.~3 \\
8  & Directory Coherence Protocols                        & Ch.~3, Ch.~8 \\
9  & Advanced Topics in Coherence                         & Ch.~3 \\
10 & Consistency and Coherence for Heterogeneous Systems  & Ch.~12 \\
11 & Specifying and Validating Memory Consistency Models  & Ch.~10 \\
\bottomrule
\end{tabular}
\caption{Chapter listing of the Primer, second edition. Page ranges vary between
the Springer and OAPEN printings and should be cited from the reader's own copy.}
\label{tab:toc-primer}
\end{table}

\noindent \textbf{Use in this paper.} This is the standard reference on the
subject and the source of the SWMR and data-value invariants stated in
\S\ref{sec:invariants}. Chapters~3--5 are the direct source for the model
definitions of Chapter~4; Chapter~11 informs the verification methodology of
Chapter~10.
Where its notation and Sarangi's differ, this paper follows Sarangi's
and notes the correspondence.

\section{Tier 2: Supporting Texts}

\subsection{Hennessy and Patterson, \emph{Computer Architecture: A Quantitative
Approach}}

Sixth edition, Morgan Kaufmann, 2019.
Chapter~5, \emph{Thread-Level Parallelism},
is the relevant portion:

\begin{table}[H]
\centering
\small
\begin{tabular}{@{}ll@{}}
\toprule
\textbf{\S} & \textbf{Title} \\
\midrule
5.1 & Introduction \\
5.2 & Centralized Shared-Memory Architectures \\
5.3 & Performance of Symmetric Shared-Memory Multiprocessors \\
5.4 & Distributed Shared-Memory and Directory-Based Coherence \\
5.5 & Synchronization: The Basics \\
5.6 & \textbf{Models of Memory Consistency: An Introduction} \\
5.7 & Cross-Cutting Issues \\
5.8 & Putting It All Together: Multicore Processors and Their Performance \\
\bottomrule
\end{tabular}
\caption{Hennessy \& Patterson, 6th ed., Chapter 5.}
\label{tab:toc-hp}
\end{table}

\noindent \textbf{Use in this paper.} \S5.6 is a compact treatment of consistency
that is useful as a cross-check but does not go beyond what the Tier~1 sources
provide.
\S5.3 and \S5.8 are consulted for the benchmarking methodology of
Chapter~6, since the quantitative framework of this text is the standard one for
reporting multiprocessor performance.

\subsection{Culler, Singh and Gupta, \emph{Parallel Computer Architecture}}

Morgan Kaufmann, 1998.
Dated in its hardware examples, but it contains the most
extended textbook treatment of relaxed consistency models available, written
while the debate this paper examines was still live --- which makes it a primary
source for Chapter~5 as much as a textbook.

\begin{table}[H]
\centering
\small
\begin{tabular}{@{}rll@{}}
\toprule
\textbf{Ch.} & \textbf{Title} & \textbf{Used in} \\
\midrule
5 & Shared Memory Multiprocessors (\S5.2 coherence, \S5.3 consistency) & Ch.~3, Ch.~4 \\
6 & Snoop-based Multiprocessor Design                                  & Ch.~3 \\
8 & Directory-based Cache Coherence                                    & Ch.~3 \\
9 & Hardware-Software Tradeoffs (\S9.1 relaxed memory consistency models) & Ch.~4, Ch.~5 \\
\bottomrule
\end{tabular}
\caption{Culler, Singh \& Gupta, chapters consulted.}
\label{tab:toc-culler}
\end{table}

\subsection{Herlihy, Shavit, Luchangco and Spear, \emph{The Art of
Multiprocessor Programming}}

Second edition, Morgan Kaufmann, 2020.
Consulted for the software side: the
correctness conditions for concurrent objects, spin-lock implementations, and the
lock-free structures used as simulator workloads in Chapter~6
(Treiber stack, Michael--Scott queue).
Appendix~B of that text covers the hardware
memory model from the programmer's perspective and is a useful counterweight to
the architect's framing of the Tier~1 sources.

\section{Tier 3: Architecture Manuals}

Where this paper states what a particular instruction set guarantees, the vendor
manual is the authoritative source and is cited in preference to any textbook.

\begin{table}[H]
\centering
\small
\begin{tabular}{@{}p{3.1cm}p{5.0cm}p{5.2cm}@{}}
\toprule
\textbf{ISA} & \textbf{Document} & \textbf{Relevant portion} \\
\midrule
x86-64 & Intel 64 and IA-32 Architectures Software Developer's Manual, Vol.~3A
       & Chapter on Memory Ordering; the store-buffer and \code{MFENCE}
         specification \\
AArch64 & Arm Architecture Reference Manual for A-profile
        & Ch.~B2, The AArch64 Application Level Memory Model; \code{DMB},
          \code{LDAR}, \code{STLR} \\
RISC-V & The RISC-V Instruction Set Manual, Vol.~I: Unprivileged ISA
       & The RVWMO chapter and the \code{FENCE} bit encoding \\
Linux  & \code{Documentation/memory-barriers.txt} and the Linux Kernel Memory
         Model (\code{tools/memory-model})
       & An operational specification maintained by practitioners; used in
         Ch.~5 and Ch.~11 \\
\bottomrule
\end{tabular}
\caption{Architecture manuals consulted.}
\label{tab:manuals}
\end{table}

\section{Research Literature}

The papers cited in Chapter~5 are the primary literature of the field and are
listed in the bibliography rather than tabulated here.
Six are load-bearing for
this paper's argument and are read in full rather than consulted:
Lamport~\cite{lamport1979multiprocessor}, Dubois et
al.~\cite{dubois1986memory}, Adve and Hill~\cite{adve1990weak}, Adve and
Gharachorloo's tutorial~\cite{adve1996shared}, Gniady et
al.~\cite{gniady1999sc}, and Alglave et al.~\cite{alglave2014herding}.

\section{A Note on Notation}

Four of the sources above use incompatible notation for the same objects.
The
conventions adopted in this paper are those of Sarangi's
\emph{Next-Gen} \S9.4, on the grounds of continuity with the course.
Table~\ref{tab:notation} records the correspondence so that a reader working from
a different source can follow.

\begin{table}[H]
\centering
\small
\begin{tabular}{@{}llll@{}}
\toprule
\textbf{This paper} & \textbf{Meaning} & \textbf{Primer} & \textbf{Alglave et al.} \\
\midrule
$\po$   & program order                & program order      & \texttt{po} \\
$\rf$   & reads-from                   & \texttt{rf}        & \texttt{rf} \\
$\co$   & coherence order              & memory order (per address) & \texttt{co} \\
$\fr$   & from-read                    & --- (derived)      & \texttt{fr} \\
$\ppo$  & preserved program order      & ordering rules     & \texttt{ppo} \\
PLSC    & per-location sequential consistency & SWMR + data value & \texttt{sc-per-location} \\
\bottomrule
\end{tabular}
\caption{Correspondence of notation across sources.}
\label{tab:notation}
\end{table}
